\chapter{Dudko--Lyubich and current MLC cases}

\section{A common theorem shape}

The recent literature proves local connectivity for different
renormalization classes.  Their analytic arguments are not identical, but
the formal end point is uniform:
\[
 \text{class predicate}
\ +\ \text{connected shrinking parameter pieces}
\ \Longrightarrow\ \text{MLC at }c.
\]
The Lean structure \lean{LiteratureMLCData} records this common shape.  Its
fields are:
\begin{enumerate}
\item membership in the source-specific parameter class;
\item a connected shrinking \lean{ParameterPuzzleBasis};
\item a source key identifying the analytic theorem that produced the data.
\end{enumerate}
The theorem
\lean{localConnectivity_of_literatureMLCData} is proved once and reused by
the source-specific certificates.

\section{Bounded-type satellite parameters}

\begin{definition}[Bounded-type satellite class]
\label{def:bounded-satellite}
A bounded-type satellite predicate is a predicate
$\mathsf{Sat}_{\mathrm{bt}}:\C\to\Prop$ which records an infinite
satellite renormalization tower together with a uniform bound on its
combinatorial rotation data.
\end{definition}

\begin{theorem}[Dudko--Lyubich satellite MLC]
\label{thm:dudko-lyubich-satellite}
For the bounded-type satellite class covered by Dudko--Lyubich, the
Mandelbrot set is locally connected at every parameter in the class.
\end{theorem}

The formal certificate is
\lean{BoundedTypeSatelliteCertificate}.  Its class predicate is a parameter,
because the combinatorial definition of bounded type has to be built from
the renormalization tower rather than guessed from a period sequence.
Once the analytic proof constructs the parameter-puzzle basis, the theorem
\lean{localConnectivity_of_boundedTypeSatelliteCertificate} gives the
topological conclusion.

\section{Feigenbaum points}

\begin{definition}[Feigenbaum class]
\label{def:feigenbaum-class}
A Feigenbaum predicate $\mathsf{Fbg}:\C\to\Prop$ records the primitive
period-doubling (or the corresponding bounded combinatorial) renormalization
data together with the limiting parameter.
\end{definition}

\begin{theorem}[Dudko--Lyubich Feigenbaum MLC]
\label{thm:dudko-lyubich-feigenbaum}
For the Feigenbaum points covered by Dudko--Lyubich, the associated
principal-nest bounds and parapuzzle construction imply local connectivity
of $\M$ at the parameter.
\end{theorem}

The Lean endpoint is
\lean{FeigenbaumCertificate}.  The certificate deliberately does not
identify ``Feigenbaum'' with an informal numerical address: the exact
combinatorics must be supplied as the predicate $\mathsf{Fbg}$.

\section{Parabolically bounded primitive renormalization}

Kahn--Kapiamba--Lyubich prove MLC for a class of infinitely renormalizable
primitive parameters which may approach the cusp.  Their geometric tools
include:
\begin{itemize}
\item the Thin--Thick Decomposition;
\item Value Calculus;
\item the Wanderers Theorem;
\item the Wave Lemma.
\end{itemize}
These estimates control degeneration that is not covered by a fixed
compactness assumption on the renormalization types.

\begin{theorem}[Parabolically bounded primitive MLC]
\label{thm:parabolic-primitive-mlc}
For the parabolically bounded primitive class of
\lean{arXiv:2606.27272}, the Mandelbrot set is locally connected at every
parameter covered by the theorem.
\end{theorem}

The formal endpoint is
\lean{ParabolicallyBoundedPrimitiveCertificate}.  The missing work is the
formal definition of the class and the formalization of the four geometric
tools that construct its parameter-puzzle basis.

\section{Elephant-eye a priori bounds}

Kahn--Lyubich's
\lean{arXiv:2601.21905} develops a priori bounds for another family of
infinitely renormalizable quadratics whose little Mandelbrot copies can
approach a cusp.  This paper is not entered in the ledger as a standalone
MLC theorem unless its exact parameter-plane conclusion is extracted and
formalized.  Its role in the current program is more precise:
\[
\begin{gathered}
\text{bordered-surface geometry}\\
\Longrightarrow\ \text{uniform thin--thick control}\\
\Longrightarrow\ \text{candidate principal-nest bounds}.
\end{gathered}
\]
The resulting analytic data can feed the same certificate
\begin{center}
\lean{ParabolicallyBoundedPrimitiveCertificate}
\end{center}
or a future
source-specific certificate, but the class predicate and parameter-puzzle
construction must remain explicit.

\section{Neutral renormalization}

The neutral regime is not an MLC theorem by itself.  The recent papers
\lean{arXiv:2609.15500} and \lean{arXiv:2609.37383} provide machinery:
bubble bounds, pseudo-bubbles, neutral cascades, butterfly bounds, and
rigidity of neutral towers.  The analytic role of these results is:
\[
\text{neutral geometry}
\Longrightarrow
\text{controlled renormalization}
\Longrightarrow
\text{parameter-puzzle data}.
\]
The hybrid-class theorem \lean{arXiv:2610.25983} supplies an analytic
parameter-space layer for generalized polynomial-like maps, but it does not
by itself assert MLC.

\begin{openestimate}[Neutral-to-parameter bridge]
\label{est:neutral-parameter}
Formalize bubble and neutral-cascade geometry, state the rigidity theorem for
marked towers, and construct the connected shrinking parameter pieces needed
by \lean{NeutralRenormalizationCertificate}.
\end{openestimate}

\section{Status of the current contracts}

\begin{center}\small
\begin{tabular}{@{}P{0.23\textwidth}@{\hspace{4pt}}P{0.58\textwidth}@{\hspace{4pt}}P{0.12\textwidth}@{}}
\textbf{Class} & \textbf{Lean contract} & \textbf{Status}\\
\hline
Bounded-type satellite &
\lean{BoundedTypeSatelliteCertificate} &
interface\\
Feigenbaum &
\lean{FeigenbaumCertificate} &
interface\\
Parabolic primitive &
\lean{ParabolicallyBoundedPrimitiveCertificate} &
interface\\
Neutral renormalization &
\lean{NeutralRenormalizationCertificate} &
machinery open\\
\end{tabular}
\end{center}

The certificates prevent a common formalization error: proving a generic
topological lemma and then silently treating the published analytic estimate
as if it had already been verified in Lean.
